\chapter{Plateformes multigérants}\label{ch:in:multi-manager-platforms}

En janvier 2016, le conseiller enregistré de la plus grande plateforme multigérants déclarait 1\,625 salariés à la SEC~; dix ans
plus tard, il en déclarait 6\,140, et sa page d'accueil recensait plus de 360 équipes d'investissement. Une deuxième plateforme est passée de 287
salariés à 2\,063 sur la même décennie. Aucun pan du secteur n'a recruté aussi vite. Pour les personnes qu'elle recrute --- gérants de
portefeuille, analystes, chercheurs, ainsi que les ingénieurs et les responsables des risques qui les entourent --- une plateforme est un emploi dont la sortie est connue
d'avance~: une équipe qui perd plus qu'une fraction fixée du capital qui lui a été confié est réduite, puis fermée. Ce
chapitre décrit les plateformes en tant qu'employeurs, lit leur croissance dans leurs déclarations et mesure ce qu'une règle de drawdown
signifie pour la durée de vie d'une équipe.

\section{La plateforme vue du côté du salarié}

Le livre~16 (chapitre~3) décrit l'économie de la plateforme~: des pods qui gèrent chacun une part du capital d'un même fonds, des coûts répercutés
sur les investisseurs, des rémunérations assises sur le profit propre de chaque pod, un book central qui couvre ce que les pods détiennent en commun. Du côté
du salarié, la même structure se présente ainsi.
\begin{itemize}
\item \textbf{Une équipe est une petite entreprise au sein d'une grande.} Un gérant de portefeuille recrute des analystes et un
ou deux chercheurs, convient d'une allocation de capital et de limites de risque, et touche une part du profit propre de l'équipe, après déduction de ses coûts.
\item \textbf{Tout le reste est fourni.} Les données, l'exécution, la technologie, les systèmes de risque, la conformité, les opérations et
le financement relèvent de la plateforme~; une équipe les utilise et les paie par ses coûts.
\item \textbf{Les limites tiennent lieu de contrat.} Le drawdown d'une équipe depuis son plus haut, mesuré sur le capital qui lui a été alloué,
déclenche une réduction du capital, puis la fermeture de l'équipe (l'échelle de réduction du risque du livre~16, chapitre~8).
\item \textbf{La plateforme emploie beaucoup de personnes hors des équipes}~: le personnel des risques, de la technologie et des données qui fait tourner la
mécanique commune, et les traders du book central.
\end{itemize}

Les plateformes citées dans ce chapitre se décrivent dans ces termes~: une société qui affiche une «~approche multistratégies, multi-PM~»
dont «~165 équipes d'investissement couvrent six stratégies~»~; une autre qui est passée «~d'une société de trading pour compte propre à fortes racines systématiques
à un hedge fund multistratégies et multigérants~»~; «~une entreprise
d'investissement mondiale, multistratégies et multigérants~».

\section{Taille~: encours, salariés et équipes}

Les conseillers américains des plateformes déposent un formulaire ADV (chapitre~4), et les fichiers passés de la SEC conservent ces réponses, de sorte que la croissance d'une plateforme
peut se lire dans ses propres déclarations.

\begin{dated}{2026-01}{Quatre plateformes dans leurs déclarations et sur leurs pages}\label{dat:in:multi-manager-platforms:size}
\footnotesize
\begin{tabular}{@{}lrrrrr@{}}
\toprule
conseiller enregistré & janv. 2016 & janv. 2017 & déc. 2022 & déc. 2024 & janv. 2026\\
\midrule
Millennium Management, salariés & 1\,625 & 1\,830 & 3\,985 & 5\,550 & 6\,140\\
\quad RAUM (milliards de \$) & 181.5 & 207.9 & 341.0 & 505.9 & 571.1\\
Balyasny Asset Management, salariés & 287 & 395 & 1\,048 & 1\,857 & 2\,063\\
\quad RAUM (milliards de \$) & 29.0 & 36.0 & 181.0 & 248.0 & 265.2\\
Schonfeld Strategic Advisors, salariés & -- & 89 & 620 & 952 & 872\\
ExodusPoint Capital Management, salariés & -- & -- & 688 & 647 & 634\\
\bottomrule
\end{tabular}

\smallskip\noindent Les pages des sociétés elles-mêmes~: Millennium, «~plus de 7\,000 salariés dans le monde~», «~plus de 360 équipes d'investissement~», «~plus de 97~Md\$
d'AUM~»~; Balyasny, «~165 équipes d'investissement~» réparties en six stratégies, «~plus de 2\,000 professionnels de l'investissement et du support~»,
«~37~Md\$ d'encours sous gestion~». Le RAUM est brut (chapitre~4)~; les chiffres d'AUM des sociétés sont nets.
\end{dated}

Les déclarations disent trois choses que les pages des sociétés ne disent pas
(\cref{fig:in:multi-manager-platforms:growth}). La croissance a été rapide mais inégale~: les effectifs du plus grand conseiller ont été multipliés par 3.8
en dix ans, ceux d'un second par 7.2~; une plateforme plus jeune a vu ses effectifs baisser légèrement après 2022. Les encours ont grandi à peu près aussi vite
que les effectifs~: le RAUM par salarié du plus grand conseiller a baissé de 17~\% sur la décennie, celui du second a augmenté de 27~\%. Et la part des salariés
affectés aux fonctions de conseil a reculé chez les deux, de 63~\% à 47~\% et de 64~\% à 42~\%~: les plateformes ont construit leur mécanique commune
plus vite que leurs équipes. Pour un candidat, c'est la forme du marché de l'emploi à l'intérieur d'une plateforme~: une part plus grande des
nouveaux postes se trouve dans la technologie, les données et les risques que dans les équipes.

\begin{omfigure}
\begin{tikzpicture}
\begin{axis}[width=11cm,height=5.6cm,xmin=2015.5,xmax=2026.5,ymin=0,ymax=7,xlabel={\small date du fichier SEC},
  ylabel={\small salariés (milliers)},xtick={2016,2018,2020,2022,2024,2026},xticklabel style={/pgf/number format/1000 sep={}},
  tick label style={font=\footnotesize},grid=major,grid style={gray!20},table/col sep=comma,
  legend style={font=\scriptsize,at={(0.5,-0.3)},anchor=north,legend cell align=left,/tikz/every even column/.append style={column sep=8pt}},legend columns=2]
\addplot[omDef,thick,mark=*] table[x=t,y=employees]{figdata/industry/05-multi-manager-platforms/adv_mlp.csv};
\addplot[omThm,thick,mark=square*] table[x=t,y=employees]{figdata/industry/05-multi-manager-platforms/adv_bam.csv};
\addplot[omProp,thick,mark=triangle*] table[x=t,y=employees]{figdata/industry/05-multi-manager-platforms/adv_sch.csv};
\addplot[gray,thick,mark=diamond*] table[x=t,y=employees]{figdata/industry/05-multi-manager-platforms/adv_exo.csv};
\legend{Millennium Management,Balyasny Asset Management,Schonfeld Strategic Advisors,ExodusPoint Capital Management}
\end{axis}
\end{tikzpicture}

\omcaption{Salariés déclarés sur le formulaire ADV par les conseillers américains de quatre plateformes, dans les fichiers de la SEC de janvier 2016, janvier
2017, décembre 2022, décembre 2024 et janvier 2026 (aucun fichier entre 2017 et 2022 n'a été lu~; les traits relient les points
lus). Données~: fichiers des conseillers de la SEC, rubrique 5.A, via \texttt{in\_platforms.panel}.}\label{fig:in:multi-manager-platforms:growth}
\end{omfigure}

\section{Rotation des équipes~: ce que montrent les sources publiques}

Le nombre d'équipes qu'une plateforme recrute et ferme chaque année est le chiffre qu'un gérant de portefeuille recherche le plus, et celui qu'aucune
plateforme ne publie. La presse rapporte les arrivées et les départs d'équipes prises isolément, et parfois un total pour une année~;
aucune de ces informations n'est une source primaire, et le présent livre ne les reproduit pas. Ce que l'on peut dire à partir des sources publiques est structurel.

\begin{definition}[Ancienneté d'une équipe, taux de rotation des équipes]\label{def:in:multi-manager-platforms:tenure}
L'\emph{ancienneté}\index{ancienneté} d'une équipe est le temps écoulé entre son arrivée sur une plateforme et sa fermeture ou son départ. Le
\emph{taux de rotation des équipes}\index{taux de rotation des équipes} d'une plateforme est le nombre d'équipes qui ferment ou quittent la plateforme au cours d'une année
divisé par le nombre moyen d'équipes en activité pendant cette période.
\end{definition}

Les deux sont liées~: si les équipes quittaient la plateforme à un taux constant $\lambda$ par an, l'\omterm{def:in:multi-manager-platforms:tenure}{ancienneté} suivrait une loi exponentielle de médiane
$\ln 2/\lambda$, et une plateforme dont le taux de rotation est de 20~\% verrait la moitié des nouvelles équipes de chaque année partir dans un délai d'environ
3.5 ans. Le taux n'est pas constant en pratique~: une équipe nouvelle est la plus exposée lors de sa première année, avant que ses gains lui
donnent de la marge sous le plus haut. La section suivante calcule la courbe.

\section{Limites de drawdown et durée de vie d'une équipe}

\begin{method}[Ancienneté sous une échelle de drawdown]\label{met:in:multi-manager-platforms:tenure}
\begin{enumerate}
\item Attribuer à chaque équipe un ratio de Sharpe vrai et une volatilité de P\&L sur le capital qui lui est alloué.
\item Faire tourner son P\&L quotidien~; diviser par deux son capital lorsque son drawdown depuis le plus haut enregistré depuis son démarrage dépasse le seuil de réduction, et
la fermer lorsque le drawdown dépasse le seuil d'arrêt.
\item Noter la date de fermeture, ou l'horizon si elle survit.
\item Estimer la courbe de survie et l'\omterm{def:in:multi-manager-platforms:tenure}{ancienneté} médiane~; pour un mélange d'équipes, le taux de rotation et la part des équipes fermées
dont le talent était avéré.
\end{enumerate}
\end{method}

Les équipes du chapitre (à titre illustratif, et non selon les conditions d'une plateforme quelconque) ont une volatilité annuelle de 10~\% sur leur capital. Deux échelles sont comparées~: une échelle serrée
qui divise par deux le capital après un drawdown de 5~\% et ferme l'équipe après 7.5~\%, et une échelle souple, à 10~\% et 15~\%.
Le drawdown est mesuré depuis le démarrage de l'équipe, de sorte qu'une équipe qui a gagné de l'argent dispose d'une marge.

\begin{center}\footnotesize
\begin{tabular}{@{}>{\raggedright\arraybackslash}p{2.6cm}>{\centering\arraybackslash}p{1.9cm}>{\centering\arraybackslash}p{1.9cm}>{\centering\arraybackslash}p{1.9cm}>{\centering\arraybackslash}p{1.9cm}>{\centering\arraybackslash}p{1.9cm}@{}}
\toprule
échelle, ratio de Sharpe vrai & \omterm{def:in:multi-manager-platforms:tenure}{ancienneté} médiane & fermée en 10 ans & encore active à 1 an & à 3 ans & taux de rotation\\
\midrule
serrée, 0.0 & 0.9 an & 99.6~\% & 45.6~\% & 15.9~\% & 63~\%\\
serrée, 0.5 & 1.5 ans & 94.4~\% & 60.2~\% & 33.9~\% & 33~\%\\
serrée, 1.0 & 3.6 ans & 74.8~\% & 72.2~\% & 53.5~\% & 16~\%\\
serrée, 1.5 & plus de 10 ans & 47.9~\% & 81.8~\% & 69.5~\% & 7~\%\\
souple, 0.5 & plus de 10 ans & 41.4~\% & 96.4~\% & 79.1~\% & 6~\%\\
souple, 1.0 & plus de 10 ans & 15.3~\% & 98.7~\% & 92.8~\% & 2~\%\\
\bottomrule
\end{tabular}
\end{center}

Avec l'échelle serrée, même une bonne équipe a plus de chances d'être fermée dans les dix ans~: une équipe
de ratio de Sharpe 1.0 a une \omterm{def:in:multi-manager-platforms:tenure}{ancienneté} médiane de 3.6 ans, et une équipe de ratio de Sharpe 0.5 de 1.5 ans
(\cref{fig:in:multi-manager-platforms:surv}). L'échelle souple conserve la plupart des bonnes équipes pendant une décennie. Sur une plateforme dont les équipes sont réparties
uniformément entre des ratios de Sharpe de 0, 0.5, 1.0 et 1.5, l'échelle serrée ferme 20~\% des équipes par an, et l'échelle souple 5~\%~;
35~\% des équipes sont fermées dès la première année avec l'échelle serrée, contre 3~\% avec l'échelle souple.

\begin{omfigure}
\begin{tikzpicture}
\begin{axis}[width=11cm,height=5.8cm,xmin=0,xmax=10,ymin=0,ymax=1.02,xlabel={\small années depuis le démarrage de l'équipe},
  ylabel={\small part des équipes encore actives},tick label style={font=\footnotesize},grid=major,grid style={gray!20},
  table/col sep=comma,legend style={legend cell align=left,font=\scriptsize,at={(0.5,-0.3)},anchor=north,/tikz/every even column/.append style={column sep=8pt}},legend columns=2]
\addplot[omThm,thick] table[x=years,y=tight05]{figdata/industry/05-multi-manager-platforms/survival.csv};
\addplot[omThm,thick,dashed] table[x=years,y=tight10]{figdata/industry/05-multi-manager-platforms/survival.csv};
\addplot[omDef,thick] table[x=years,y=loose05]{figdata/industry/05-multi-manager-platforms/survival.csv};
\addplot[omDef,thick,dashed] table[x=years,y=loose10]{figdata/industry/05-multi-manager-platforms/survival.csv};
\legend{{échelle serrée, ratio de Sharpe 0.5},{échelle serrée, ratio de Sharpe 1.0},{échelle souple, ratio de Sharpe 0.5},{échelle souple, ratio de Sharpe 1.0}}
\end{axis}
\end{tikzpicture}

\omcaption{Survie des équipes simulées sous deux échelles de drawdown~: part encore active après chaque année, 4\,000 équipes
par courbe, volatilité annuelle de 10~\% du capital alloué. Échelle serrée~: capital divisé par deux à un drawdown de 5~\% depuis le plus haut atteint depuis
le démarrage, équipe fermée à 7.5~\%~; échelle souple~: 10~\% et 15~\%. Paramètres illustratifs. Données~:
\texttt{firm.teamtenure.stop\_times}, via \texttt{in\_platforms.times}.}\label{fig:in:multi-manager-platforms:surv}
\end{omfigure}

L'échelle ne permet pas de séparer vite le talent de la chance. Parmi les équipes fermées par l'échelle serrée, 39~\% avaient un ratio de Sharpe vrai
de 1.0 ou plus~; avec l'échelle souple, 14~\%. L'étape de division par deux compte~: sans elle, une équipe de ratio de Sharpe 0.5 vit en médiane 0.7 an sous le seuil d'arrêt serré,
au lieu de 1.5 ans, car la coupure divise par deux l'ampleur des pertes qui suivent un mauvais départ.
Une plateforme choisit où se placer entre fermer vite les mauvaises équipes et garder les bonnes, et un gérant de portefeuille
qui la rejoint doit demander où elle se situe.

\begin{remark}[Ce que le modèle laisse de côté]\label{rem:in:multi-manager-platforms:model}
Les plateformes réelles évaluent les équipes sur davantage que les drawdowns~: rendements rapportés au risque alloué, encombrement (crowding) avec d'autres équipes,
comportement face aux limites. Les équipes partent aussi d'elles-mêmes, vers une autre plateforme ou pour lancer un fonds (livre~16, chapitre~25).
Le capital est réalloué à la hausse après les gains. Le modèle isole la seule règle à laquelle une équipe souscrit~; il n'est pas une prévision
du taux de rotation d'une plateforme quelconque, et les volatilités et ratios de Sharpe des équipes sont choisis à titre illustratif.
\end{remark}

\section{Comment les plateformes recrutent}

Une plateforme recrute un gérant de portefeuille sur un historique et un plan. L'historique est le P\&L passé du gérant et le risque pris pour l'obtenir,
que le livre~16 (chapitre~25) juge transférable seulement s'il peut être attribué et vérifié~; le plan est le capital
que le gérant demande, la stratégie, l'équipe et le budget. Les analystes et les chercheurs sont recrutés par les gérants
de portefeuille, souvent parmi les équipes d'autres plateformes ou dans les desks des banques~; les ingénieurs et le personnel de risque sont recrutés par la plateforme elle-même.
Comme l'économie de chaque équipe lui est propre, le recrutement d'une plateforme fonctionne comme un marché~: les équipes se disputent le capital et les talents,
et une équipe fermée libère les deux.

Pour un analyste, la conséquence est que l'employeur est, en pratique, l'équipe~: quand celle-ci ferme, le poste
prend en général fin avec elle, sauf si une autre équipe recrute l'analyste. L'\omterm{def:in:multi-manager-platforms:tenure}{ancienneté} sous une échelle est donc aussi la sécurité d'emploi de l'analyste,
et le tableau ci-dessus est une raison de s'en renseigner avant de rejoindre une équipe.

\section{Tutoriel~: croissance des dépôts et ancienneté dans le modèle}\label{tut:in:multi-manager-platforms:tut}

\begin{tutorial}
\textbf{Objectif.} Lire la croissance de quatre plateformes à partir de leurs dépôts Form ADV, et calculer l'\omterm{def:in:multi-manager-platforms:tenure}{ancienneté} des équipes sous deux échelles.
\textbf{État final~:} l'encadré daté, la \cref{fig:in:multi-manager-platforms:growth}, le tableau d'\omterm{def:in:multi-manager-platforms:tenure}{ancienneté} et
la \cref{fig:in:multi-manager-platforms:surv}.
\begin{tutsteps}
\item \textbf{Le panel.} \texttt{in\_platform\_derive.py} lit cinq fichiers de conseillers de la SEC (de 2016 à 2026, aux formats xlsx et
CSV) avec \texttt{firm.formadv.read(path, hedge\_only=False)} et conserve les quatre conseillers grâce à leurs numéros CRD~;
\texttt{in\_platforms.growth} compare deux dates.
\item \textbf{Une équipe.} \texttt{firm.teamtenure.stop\_times} fait tourner un lot d'équipes sous une échelle
(voir le \cref{lst:in:multi-manager-platforms:stop})~; les rendements viennent de \texttt{firm.podshop.pods} du livre~16.
\omcode{code/firm/teamtenure/firm_teamtenure.py}{35}{48}{L'échelle appliquée jour par jour~: division par deux au seuil de réduction, fermeture au
seuil d'arrêt, les deux mesurées depuis le plus haut atteint depuis le démarrage de l'équipe.}\label{lst:in:multi-manager-platforms:stop}
\item \textbf{Le tableau.} \texttt{in\_platforms.table()} donne l'\omterm{def:in:multi-manager-platforms:tenure}{ancienneté} médiane, la part fermée, la survie à un et
trois ans et le taux de rotation pour chaque échelle et chaque ratio de Sharpe~; \texttt{mixed()} donne les chiffres au niveau de la plateforme.
\item \textbf{Vérifier le modèle.} Faire tourner l'échelle serrée avec quatre fois plus de pas par an~: l'\omterm{def:in:multi-manager-platforms:tenure}{ancienneté} médiane d'une équipe
de ratio de Sharpe 0.5 reste à moins de 10~\% de 1.5 ans. Supprimer l'étape de division par deux~: elle tombe à 0.7 an.
\end{tutsteps}
\end{tutorial}

\textbf{Pour aller plus loin.} Porter la volatilité des équipes à 15~\% et observer comment le taux de rotation de l'échelle serrée change
(exercice~7)~; donner aux équipes un ratio de Sharpe qui tombe à zéro après deux ans et demander à quelle vitesse chaque échelle les ferme.

\section{Construction~: le modèle d'ancienneté}\label{bld:in:multi-manager-platforms:teamtenure}

\begin{build}
\bfield{Rôle} Mesurer l'effet d'une règle de drawdown sur la durée de vie des équipes, pour ce chapitre et pour
l'accord du chapitre~22 avec le gérant de portefeuille.
\bfield{Interface} \texttt{firm.teamtenure}~: \texttt{Ladder(cut, stop)}~;
\texttt{stop\_times(sr, vol, ladder, n, years, rng, days)}~; \texttt{survival(times, grid)}~; \texttt{median\_tenure(times)}~; \texttt{turnover(times, years)}~;
\texttt{false\_cut\_share(times\_by\_sr, skilled)}. Encapsule \texttt{firm.podshop.pods}.
\bfield{Règles} Les drawdowns sont mesurés sur le capital effectivement géré, depuis le plus haut atteint depuis le démarrage~; une équipe fermée reste
fermée~; les équipes encore actives à l'horizon sont censurées et non fermées.
\bfield{Tests d'acceptation} \texttt{code/firm/teamtenure/tests/}~: survie, médiane et taux de rotation sur des durées connues~; davantage de talent
et une échelle plus souple signifient moins de fermetures~; la part de fausses coupes sur des durées construites.
\bfield{Pour aller plus loin} Réallocation du capital après les gains~; équipes qui partent d'elles-mêmes à un taux constant~; un ratio de Sharpe qui
décroît avec l'âge de l'équipe.
\end{build}

\begin{omsources}
\item SEC, Information about registered investment advisers: fichiers mensuels de janvier 2016, janvier 2017, décembre 2022,
décembre 2024 et janvier 2026 (copies de l'Internet Archive).
\item Les pages propres des plateformes~: Millennium, Balyasny, Schonfeld, ExodusPoint (registre F2--F5).
\item Livre~16, chapitres 3, 8 et 25, pour l'économie de la plateforme, l'échelle de réduction du risque (de-risking) et les historiques de performance.
\end{omsources}

\section{Exercices}

\begin{exercise}[$\star$]\label{exo:in:multi-manager-platforms:1}
À partir de l'encadré daté, calculer la croissance du nombre de salariés et des encours sous gestion du plus grand conseiller entre janvier 2016 et janvier 2026.
\end{exercise}

\begin{exercise}[$\star$]\label{exo:in:multi-manager-platforms:2}
Si les équipes quittent une plateforme à un taux constant de 20~\% par an, quelle est l'\omterm{def:in:multi-manager-platforms:tenure}{ancienneté} médiane d'une équipe~?
\end{exercise}

\begin{exercise}[$\star$]\label{exo:in:multi-manager-platforms:3}
La page d'une plateforme recense 360 équipes et 7\,000 salariés. Que cela suggère-t-il de la part du personnel en dehors des
équipes, si une équipe compte en moyenne cinq personnes~?
\end{exercise}

\begin{exercise}[$\star\star$]\label{exo:in:multi-manager-platforms:4}
Calculer la part de conseil du conseiller de la seconde plateforme en janvier 2016 et en janvier 2026, puis dire ce que la variation
signifie pour un candidat.
\end{exercise}

\begin{exercise}[$\star\star$]\label{exo:in:multi-manager-platforms:5}
Sous l'échelle serrée, quelle part des équipes dont le ratio de Sharpe vaut 1.0 est encore en activité après trois ans~? Au bout de
combien de temps la moitié d'entre elles est-elle fermée~?
\end{exercise}

\begin{exercise}[$\star\star$]\label{exo:in:multi-manager-platforms:6}
Lire la \cref{fig:in:multi-manager-platforms:surv}~: à partir de quel âge l'écart entre les échelles serrée et souple, pour les équipes
dont le ratio de Sharpe vaut 0.5, dépasse-t-il 40 points de pourcentage~?
\end{exercise}

\begin{exercise}[$\star\star\star$]\label{exo:in:multi-manager-platforms:7}
\emph{Programmation.} Relancer \texttt{stop\_times} pour un ratio de Sharpe de 1.0 sous l'échelle serrée, avec une volatilité de 15~\%
au lieu de 10~\%. Que se passe-t-il pour l'\omterm{def:in:multi-manager-platforms:tenure}{ancienneté} médiane, et pourquoi~?
\end{exercise}

\begin{exercise}[$\star\star\star$]\label{exo:in:multi-manager-platforms:8}
\emph{Trouver la faille.} «~Sous l'échelle serrée, 39~\% des équipes fermées étaient compétentes, donc la plateforme devrait cesser de fermer
les équipes.~»
\end{exercise}

\section{Problème~: Combien de temps une équipe dure-t-elle~?}

\begin{problem}[{Problème du week-end --- combien de temps une équipe dure-t-elle~?}]\label{pb:in:multi-manager-platforms:1}
Une gérante de portefeuille se voit proposer une place sur une plateforme dotée d'une échelle de drawdown. Elle estime que le ratio de Sharpe de sa stratégie se situe
entre 0.5 et 1.0 et veut connaître ses chances.

\medskip\noindent\textbf{Partie I --- La plateforme.}
\begin{enumerate}
\item Énumérer ce qu'apporte une équipe et ce qu'apporte la plateforme.
\item Donner le nombre de salariés du plus grand conseiller en janvier 2016 et en janvier 2026, ainsi que le facteur de croissance.
\item Donner la variation de ses encours sous gestion par salarié et de sa part de conseil sur la décennie.
\item Que disent les pages des deux plateformes au sujet de leur nombre d'équipes~?
\item Pourquoi le chapitre n'affiche-t-il pas les décomptes de fermetures d'équipes publiés par la presse~?
\end{enumerate}

\medskip\noindent\textbf{Partie II --- L'\omterm{def:in:multi-manager-platforms:tenure}{ancienneté}.}
\begin{enumerate}[resume]
\item Définir l'\omterm{def:in:multi-manager-platforms:tenure}{ancienneté} d'une équipe et le taux de rotation, puis relier l'\omterm{def:in:multi-manager-platforms:tenure}{ancienneté} médiane à un taux constant $\lambda$.
\item Énoncer les deux échelles et la volatilité des équipes.
\item Donner l'\omterm{def:in:multi-manager-platforms:tenure}{ancienneté} médiane pour des ratios de Sharpe de 0.5 et de 1.0 sous l'échelle serrée.
\item Donner la part de ces équipes fermées dans les dix ans, sous chacune des échelles.
\item Pourquoi une équipe nouvelle est-elle la plus exposée au cours de sa première année~?
\end{enumerate}

\medskip\noindent\textbf{Partie III --- Le point de vue de la plateforme.}
\begin{enumerate}[resume]
\item Donner le taux de rotation de la plateforme mixte sous chacune des échelles.
\item Donner la part des équipes fermées dès leur première année, sous chacune des échelles.
\item Donner la part des équipes fermées qui étaient compétentes (ratio de Sharpe de 1.0 ou plus), sous chacune des échelles.
\item Que fait l'étape de division par deux à l'\omterm{def:in:multi-manager-platforms:tenure}{ancienneté} d'une équipe dont le ratio de Sharpe vaut 0.5~?
\item Vérifier le modèle~: que devient cette médiane lorsque l'étape est divisée par quatre~?
\end{enumerate}

\medskip\noindent\textbf{Partie IV --- Le verdict.}
\begin{enumerate}[resume]
\item Énoncer le \emph{résultat nommé}~: l'\omterm{def:in:multi-manager-platforms:tenure}{ancienneté} médiane sous l'échelle serrée pour des ratios de Sharpe de 0.5 et de 1.0, et la
part des équipes fermées qui étaient compétentes sous chacune des échelles.
\item Quelles trois choses devrait-elle demander à la plateforme avant d'y entrer~?
\item Pourquoi l'\omterm{def:in:multi-manager-platforms:tenure}{ancienneté} est-elle aussi la sécurité d'emploi de ses analystes~?
\item Que laisse de côté le modèle qui pourrait allonger ou raccourcir son \omterm{def:in:multi-manager-platforms:tenure}{ancienneté}~?
\item En deux phrases, que devrait-elle conclure~?
\end{enumerate}
\end{problem}

\section{Questions d'entretien}

\begin{interviewq}[$\star$ \iqroles{trader, chercheur}]\label{iq:in:multi-manager-platforms:1}
Qu'est-ce qu'une plateforme multigérants, et en quoi le travail dans l'une de ses équipes diffère-t-il
de celui dans un fonds à stratégie unique~?
\end{interviewq}

\begin{interviewq}[$\star$ \iqroles{chercheur, risque}]\label{iq:in:multi-manager-platforms:2}
Une équipe gère une volatilité annuelle de 10~\% avec un ratio de Sharpe de 1. Quelle est la probabilité qu'elle soit en baisse de 5~\% ou plus après
sa première année~?
\end{interviewq}

\begin{interviewq}[$\star\star$ \iqroles{risque}]\label{iq:in:multi-manager-platforms:3}
Pourquoi une plateforme réduit-elle de moitié le capital d'une équipe après un premier drawdown plutôt que de la fermer aussitôt~?
\end{interviewq}

\begin{interviewq}[$\star\star$ \iqroles{chercheur}]\label{iq:in:multi-manager-platforms:4}
Une plateforme ferme chaque année 20~\% de ses équipes. Que peut-on en déduire sur la compétence des équipes qui restent~?
\end{interviewq}

\begin{interviewq}[$\star\star$ \iqroles{trader}]\label{iq:in:multi-manager-platforms:5}
En tant que nouveau gérant de portefeuille, comment dimensionneriez-vous votre risque pendant vos premiers mois sous une échelle de drawdown, et pourquoi~?
\end{interviewq}

\begin{interviewq}[$\star\star\star$ \iqroles{chercheur, risque}]\label{iq:in:multi-manager-platforms:6}
Le P\&L d'une équipe est un mouvement brownien avec dérive $\mu>0$ et volatilité $\sigma$. Montrez qu'un arrêt à une perte $d$ sous son
capital de départ l'atteint avec une probabilité $e^{-2\mu d/\sigma^2}$, tandis qu'un arrêt à $d$ sous son plus haut courant l'atteint
avec une probabilité égale à un. Que cela implique-t-il pour une échelle mesurée depuis le plus haut~?
\end{interviewq}
